% ============================================================
% M3 S3 练习件·W1 臂A —— 课时02 倾斜角与斜率（2.2.1上）＝照抄母版件（课时01）体例
% 生成：M3 成卷轮 S3 W1 臂A｜2026-09-14｜编译 cwd＝本目录（中文路径禁 \input 绝对路径）
%   xelatex main-true.tex ×2 → 含详解印本；xelatex main-false.tex ×2 → 纯题版（单源双本）
% 输入链（全只读）：题面库 课时02-倾斜角与斜率.md（题面逐字装配）＋ -答案侧.md（值逐字回嵌）
%   ＋ 定稿/批A-课时02-倾斜角与斜率.md（详解回嵌源）；toolchain＝qp-m3.sty（钉后版
%   md5 7c3930362be8a0a2bdf21bbf8ac16573，本地挂载副本，破损即停）。
% 母版体例：详见《练习件母版体例说明.md》（课时01 件）；门谱读数见 过程件 工作区/_tmpM3S3W1臂A0914/。
% 键账：件内 ansblock 键集 ≡ 题面库片练键集 ≡ manifest 练键序 E1~E16（16 键，零缺漏零多余）；
%   拓区隔离：2章-拓/测/滚 键不入本件（S4 拓展册收，本片 T1~T4 不入）；导学键 G1~G5 属导学件域，亦不入。
% 卷式例外案B：练习件不属卷式——答案块物理落点恒为题后 ansblock，禁卷末附卷制；
%   全线括线模（导言 \ansblockgrayfalse 一次置定，件内不切换；灰底盒不可跨栏断，练习件禁用）。
% 题序＝manifest 键序 E1~E16（零重排）；印面题号 1~16 连号（\tihao 号＝\ansitem 号同键）；
%   三档切位 1~7／8~14／15~16（照库序切位，非难度重排）。
% 槽配实录：单选8（E1~E5/E7~E9）＋多选1（E6，≤4 合规）＋填空3（E10~E12）＋解答4（E13~E16）；
%   10选4填2解 为波次方案基线槽配，本片实配照题面侧头标（批A §四 满足度表同口径）。
% 选项槽位：默认四连排（0.25\linewidth−0.25\lxhang），超宽降档梯 四连排→两连排
%   （0.5\linewidth−0.5\lxhang−0.25em）→单列 \lxopt；禁缩字号/负 kern/删标点凑宽。
%   本片实测降档：见值台账「降档登记」＋回执（槽宽门真算读数）。
% 解答书写区：E13/E14/E15 无小问默认 16mm；E16 四小问 32mm（16~40mm 带内；纯题档吞
%   ansblock 不吞 \liubai）。
% 填空印答：E10/E11/E12 空位 \kongda{值}（详解版印答、纯题版退定宽空线，两档等形）。
% 图债：E1 国旗五星坐标图（题面「如图」占位保留，解算不依赖图；图工位补图后回冲印面）。
% 难度词：照题面侧头标第 4 字段简/中档词逐字；知识点 N 系预排（本课时知识点一＝直线的
%   倾斜角、二＝直线的斜率、三＝倾斜角与斜率互求，照批A 课时切分；导学件课时02 尚未落盘，
%   落盘后如异动按变更协议回核）。
% ============================================================
\documentclass[fontset=none]{ctexart}
\usepackage{qp-m3}
% —— 印面逐字符号路由补钉（承导学件母版；− 走 TNR、▱ NSC 子块；禁改 sty 保 md5） ——
\xeCJKDeclareCharClass{Default}{"2212}%
\setCJKfamilyfont{nscblk}[Path=C:/提示词/工作区/字替对照-0909/variantF/fonts/,BoldFont=NSC-w600.ttf]{NSC-w500.ttf}%
\xeCJKDeclareSubCJKBlock{nscblk}{"25B1}%
\newcommand{\pxparallelogram}{{\CJKfamily{nscblk}▱}}%
% —— 断行微调（承导学件母版实测档，三零门承重；\emergencystretch 不另设） ——
\xeCJKsetup{CJKglue={\hskip 0.10em plus 0.08em minus 0.05em}}%
% —— 练习件全线括线模（S3 波次方案 §四.3：导言一次置定、件内不切换） ——
\ansblockgrayfalse
% —— 页脚件名（qp-headfoot 复用入口） ——
\renewcommand{\qpzhangming}{第二章\quad 平面解析几何}%
\renewcommand{\qpjianming}{练习件}%

\begin{document}

\keshi{第2课时\quad 倾斜角与斜率}
\begin{multicols}{2}
\emergencystretch=1em

% ============================================================
% 基础巩固（题1~7＝E1~E7：选择7＋填空0；题后紧跟答案制，全线括线 ansblock）
% ============================================================
\dangbadge{基}{础}{巩}{固}

\tihao{1}\zu{倾斜角与斜率×单选} \tieside{中档(知识点一)}2020年12月4日，嫦娥五号探测器在月球表面第一次动态展示国旗．1949年公布的《国旗制法说明》中就五星的位置规定：大五角星有一个角尖正向上方，四颗小五角星均各有一个角尖正对大五角星的中心点．有人发现，第三颗小星的姿态与大星相近．为便于研究，如图，以大星的中心点为原点，建立直角坐标系，\(OO_1\)，\(OO_2\)，\(OO_3\)，\(OO_4\)分别是大星中心点与四颗小星中心点的联结线，\(\alpha\approx 16^\circ\)，则第三颗小星的一条边\(AB\)所在直线的倾斜角约为\nobreak(\kongwei)
\bindopt {\fontsize{10.5pt}{\qpTJgLeadLiti}\selectfont \makebox[\dimexpr0.25\linewidth-0.25\lxhang\relax][l]{A．0°}\makebox[\dimexpr0.25\linewidth-0.25\lxhang\relax][l]{B．1°}\makebox[\dimexpr0.25\linewidth-0.25\lxhang\relax][l]{C．2°}\makebox[\dimexpr0.25\linewidth-0.25\lxhang\relax][l]{D．3°}\par}

\begin{ansblock}[2章-练-课时02-E1]
% ans:2章-练-课时02-E1
\ansitem{1}{C}
\ansnote{详解}{∵\(O\)、\(O_3\)都为五角星的中心点，∴\(OO_3\)平分第三颗小星的一个角；又五角星的内角为\(36^\circ\)，可知\(\angle BAO_3=18^\circ\)．过\(O_3\)作\(x\)轴的平行线\(O_3E\)，则\(\angle OO_3E=\alpha\approx 16^\circ\)，所以直线\(AB\)的倾斜角约为\(18^\circ-16^\circ=2^\circ\)．故选：C．}
\end{ansblock}

\tihao{2}\zu{倾斜角与斜率×单选} \tieside{简单(知识点二)}已知两点\(A(1,-2)\)，\(B(2,1)\)，直线\(l\)过点\(P(0,-1)\)且与线段\(AB\)有交点，则直线\(l\)的斜率的取值范围为\nobreak(\kongwei)
\bindopt {\fontsize{10.5pt}{\qpTJgLeadLiti}\selectfont \makebox[\dimexpr0.5\linewidth-0.5\lxhang-0.25em\relax][l]{A．\([-1,1]\)}\makebox[\dimexpr0.5\linewidth-0.5\lxhang-0.25em\relax][l]{B．\((-\infty,-1]\)}\par}
\bindopt {\fontsize{10.5pt}{\qpTJgLeadLiti}\selectfont \makebox[\dimexpr0.5\linewidth-0.5\lxhang-0.25em\relax][l]{C．\((-1,1)\)}\makebox[\dimexpr0.5\linewidth-0.5\lxhang-0.25em\relax][l]{D．\([1,+\infty)\)}\par}

\begin{ansblock}[2章-练-课时02-E2]
% ans:2章-练-课时02-E2
\ansitem{2}{A}
\ansnote{详解}{直线\(PA\)的斜率为\(k_{PA}=(-2+1)/(1-0)=-1\)，直线\(PB\)的斜率为\(k_{PB}=(1+1)/(2-0)=1\)．数形结合：当过\(P\)的直线\(l\)绕\(P\)从\(PA\)旋转到\(PB\)（保持与线段\(AB\)有公共点）时，其斜率介于\(k_{PA}\)与\(k_{PB}\)之间，故\(k\in[-1,1]\)．故选：A．}
\end{ansblock}

\tihao{3}\zu{倾斜角与斜率×单选} \tieside{中档(知识点二)}直线\(l\)的倾斜角为\(\alpha\)，且\(\sin\alpha=3/5\)，则直线\(l\)的斜率是\nobreak(\kongwei)

\bindopt {\fontsize{10.5pt}{\qpTJgLeadLiti}\selectfont \makebox[\dimexpr0.5\linewidth-0.5\lxhang-0.25em\relax][l]{A．\(-4/3\)}\makebox[\dimexpr0.5\linewidth-0.5\lxhang-0.25em\relax][l]{B．\(3/4\)}\par}
\bindopt {\fontsize{10.5pt}{\qpTJgLeadLiti}\selectfont \makebox[\dimexpr0.5\linewidth-0.5\lxhang-0.25em\relax][l]{C．\(3/4\)或\(-3/4\)}\makebox[\dimexpr0.5\linewidth-0.5\lxhang-0.25em\relax][l]{D．\(4/3\)或\(-4/3\)}\par}

\begin{ansblock}[2章-练-课时02-E3]
% ans:2章-练-课时02-E3
\ansitem{3}{C}
\ansnote{详解}{\(\alpha\in[0^\circ,180^\circ)\)且\(\sin\alpha=3/5>0\)，故\(\alpha\)为锐角或钝角．当\(\alpha\)为锐角时，\(\cos\alpha=4/5\)，\(k=\tan\alpha=3/4\)；当\(\alpha\)为钝角时，\(\cos\alpha=-4/5\)，\(k=\tan\alpha=-3/4\)．故直线\(l\)的斜率是\(3/4\)或\(-3/4\)．选：C．}
\end{ansblock}

\tihao{4}\zu{倾斜角与斜率×单选} \tieside{简单(知识点三)}已知直线斜率为\(k\)，且\(-1\leqslant k\leqslant\sqrt{3}\)，那么倾斜角\(\alpha\)的取值范围是\nobreak(\kongwei)

\lxopt{A．\([0,\pi/3]\cup[\pi/2,3\pi/4)\)}

\lxopt{B．\([0,\pi/3]\cup[3\pi/4,\pi)\)}

\lxopt{C．\([0,\pi/6]\cup[\pi/2,3\pi/4)\)}

\lxopt{D．\([0,\pi/6]\cup[3\pi/4,\pi)\)}

\begin{ansblock}[2章-练-课时02-E4]
% ans:2章-练-课时02-E4
\ansitem{4}{B}
\ansnote{详解}{\(\alpha\in[0,\pi)\)，\(-1\leqslant k\leqslant\sqrt{3}\) 即 \(-1\leqslant\tan\alpha\leqslant\sqrt{3}\)．当\(0\leqslant\tan\alpha\leqslant\sqrt{3}\)时，\(\alpha\in[0,\pi/3]\)；当\(-1\leqslant\tan\alpha<0\)时，\(\alpha\in[3\pi/4,\pi)\)．所以\(\alpha\in[0,\pi/3]\cup[3\pi/4,\pi)\)．故选：B．}
\end{ansblock}

\tihao{5}\zu{倾斜角与斜率×单选} \tieside{简单(知识点三)}直线\(l\)经过点\(A(2,1)\)，\(B(3,t^{2})\)（\(-\sqrt{2}\leqslant t\leqslant\sqrt{2}\)），则直线\(l\)倾斜角的取值范围是\nobreak(\kongwei)

\lxopt{A．\([0,\pi/4]\cup[3\pi/4,\pi)\)}

\lxopt{B．\([0,\pi/3]\cup[2\pi/3,\pi)\)}

\lxopt{C．\([\pi/4,\pi/2)\cup(\pi/2,3\pi/4]\)}

\lxopt{D．\([0,\pi/4]\)}

\begin{ansblock}[2章-练-课时02-E5]
% ans:2章-练-课时02-E5
\ansitem{5}{A}
\ansnote{详解}{\(k_l=(t^{2}-1)/(3-2)=t^{2}-1\)．由\(-\sqrt{2}\leqslant t\leqslant\sqrt{2}\)得\(t^{2}\in[0,2]\)，故\(k\in[-1,1]\)，即\(\tan\theta\in[-1,1]\)（\(\theta\)为倾斜角），得\(\theta\in[0,\pi/4]\cup[3\pi/4,\pi)\)．选：A．}
\end{ansblock}

\tihao{6}\duoxuan\hspace{0.6em}\zu{倾斜角与斜率×多选} \tieside{简单(知识点二)}（多选）若经过\(A(1-a,1+a)\)和\(B(3,a)\)的直线的倾斜角为钝角，则实数\(a\)的值可能为\nobreak(\kongwei)
\bindopt {\fontsize{10.5pt}{\qpTJgLeadLiti}\selectfont \makebox[\dimexpr0.25\linewidth-0.25\lxhang\relax][l]{A．\(-2\)}\makebox[\dimexpr0.25\linewidth-0.25\lxhang\relax][l]{B．\(0\)}\makebox[\dimexpr0.25\linewidth-0.25\lxhang\relax][l]{C．\(1\)}\makebox[\dimexpr0.25\linewidth-0.25\lxhang\relax][l]{D．\(2\)}\par}

\begin{ansblock}[2章-练-课时02-E6]
% ans:2章-练-课时02-E6
\ansitem{6}{BCD}
\ansnote{详解}{倾斜角为钝角\(\Longleftrightarrow\)斜率存在且为负，即\(A\)、\(B\)横坐标不等且\(k_{AB}<0\)．\(k_{AB}=(1+a-a)/[(1-a)-3]=1/(-2-a)<0\)，得\(a+2>0\)，即\(a>-2\)．又\(a=-2\)时\(A\)、\(B\)重合？\(a=-2\)时\(A(3,-1)\)、\(B(3,-2)\)，横坐标相同、斜率不存在，倾斜角\(90^\circ\)，不在“钝角”取值内且\(a>-2\)已排除．故\(a>-2\)，选项中0、1、2满足．选：BCD．}
\end{ansblock}

\tihao{7}\zu{倾斜角与斜率×单选} \tieside{简单(知识点二)}若\(A(-2,3)\)，\(B(3,2)\)，\(C(1/2,m)\)三点共线，则实数\(m\)的值为\nobreak(\kongwei)

\bindopt {\fontsize{10.5pt}{\qpTJgLeadLiti}\selectfont \makebox[\dimexpr0.25\linewidth-0.25\lxhang\relax][l]{A．\(2\)}\makebox[\dimexpr0.25\linewidth-0.25\lxhang\relax][l]{B．\(-2\)}\makebox[\dimexpr0.25\linewidth-0.25\lxhang\relax][l]{C．\(5/2\)}\makebox[\dimexpr0.25\linewidth-0.25\lxhang\relax][l]{D．\(-1/2\)}\par}

\begin{ansblock}[2章-练-课时02-E7]
% ans:2章-练-课时02-E7
\ansitem{7}{C}
\ansnote{详解}{三点共线\(\Longleftrightarrow k_{AB}=k_{AC}\)．\(k_{AB}=(2-3)/(3+2)=-1/5\)；\(k_{AC}=(m-3)/(1/2+2)=(m-3)/(5/2)=2(m-3)/5\)．令\(2(m-3)/5=-1/5\)，得\(2(m-3)=-1\)，\(m=5/2\)．故选：C．}
\end{ansblock}

% ============================================================
% 综合提升（题8~14＝E8~E14：单选2（E8/E9）＋填空3（E10~E12）＋解答2（E13/E14））
% ============================================================
\dangbadge{综}{合}{提}{升}

\tihao{8}\zu{倾斜角与斜率×单选} \tieside{简单(知识点三)}直线\(m\)过点\(O(0,0)\)，\(A(1,\sqrt{2})\)，其倾斜角为\(\alpha\)，现将直线\(m\)绕原点\(O\)逆时针旋转得到直线\(m'\)：\(y=kx\)，若直线\(m'\)的倾斜角为\(2\alpha\)，则\(k\)的值为\nobreak(\kongwei)
\bindopt {\fontsize{10.5pt}{\qpTJgLeadLiti}\selectfont \makebox[\dimexpr0.25\linewidth-0.25\lxhang\relax][l]{A．\(2\sqrt{2}\)}\makebox[\dimexpr0.25\linewidth-0.25\lxhang\relax][l]{B．\(-2\sqrt{2}\)}\makebox[\dimexpr0.25\linewidth-0.25\lxhang\relax][l]{C．\(2\)}\makebox[\dimexpr0.25\linewidth-0.25\lxhang\relax][l]{D．\(-2\)}\par}

\begin{ansblock}[2章-练-课时02-E8]
% ans:2章-练-课时02-E8
\ansitem{8}{B}
\ansnote{详解}{\(\tan\alpha=k_{OA}=\sqrt{2}\)．\(m'\)的倾斜角为\(2\alpha\)，故 \(k=\tan 2\alpha=2\tan\alpha/(1-\tan^{2}\alpha)=2\sqrt{2}/(1-2)=-2\sqrt{2}\)．选：B．（合理性：\(\tan\alpha=\sqrt{2}>1\)，\(\alpha\in(\pi/4,\pi/2)\)，\(2\alpha\in(\pi/2,\pi)\)确为钝角，\(\tan 2\alpha<0\)）}
\end{ansblock}

\tihao{9}\zu{倾斜角与斜率×单选} \tieside{简单(知识点三)}直线\(2x\cdot\sin 210^\circ-y-2=0\)的倾斜角是\nobreak(\kongwei)
\bindopt {\fontsize{10.5pt}{\qpTJgLeadLiti}\selectfont \makebox[\dimexpr0.25\linewidth-0.25\lxhang\relax][l]{A．45°}\makebox[\dimexpr0.25\linewidth-0.25\lxhang\relax][l]{B．135°}\makebox[\dimexpr0.25\linewidth-0.25\lxhang\relax][l]{C．30°}\makebox[\dimexpr0.25\linewidth-0.25\lxhang\relax][l]{D．150°}\par}

\begin{ansblock}[2章-练-课时02-E9]
% ans:2章-练-课时02-E9
\ansitem{9}{B}
\ansnote{详解}{\(\sin 210^\circ=-1/2\)，直线方程即\(-x-y-2=0\)，\(y=-x-2\)，斜率\(k=2\sin 210^\circ=-1\)，故倾斜角为\(135^\circ\)．选：B．}
\end{ansblock}

\tihao{10}\zu{倾斜角与斜率×填空} \tieside{简单(知识点二)}若\(A(3,1)\)，\(B(-2,k)\)，\(C(8,1)\)三点能构成三角形，则实数\(k\)的取值范围为\kongda{(−∞,1)∪(1,+∞)}．

\begin{ansblock}[2章-练-课时02-E10]
% ans:2章-练-课时02-E10
\ansitem{10}{(−∞,1)∪(1,+∞)}
\ansnote{详解}{三点能构成三角形\(\Longleftrightarrow\)三点不共线．由\(k_{AB}\neq k_{AC}\)，即 \((k-1)/(-2-3)\neq(1-1)/(8-3)\)，得\((k-1)/(-5)\neq0\)，\(k-1\neq0\)，\(k\neq1\)．故\(k\)的取值范围为\((-\infty,1)\cup(1,+\infty)\)．}
\end{ansblock}

\tihao{11}\zu{倾斜角与斜率×填空} \tieside{简单(知识点三)}若直线\(l\)的倾斜角的变化范围为\([\pi/6,\pi/3)\)，则直线斜率的取值范围是\kongda{[√3/3, √3)}．

\begin{ansblock}[2章-练-课时02-E11]
% ans:2章-练-课时02-E11
\ansitem{11}{[√3/3, √3)}
\ansnote{详解}{正切函数\(k=\tan\alpha\)在\([\pi/6,\pi/3)\)上单调递增，\par\noindent 故\(\alpha\in[\pi/6,\pi/3)\)时 \(\tan\alpha\in[\tan(\pi/6),\ \tan(\pi/3))=[\sqrt{3}/3,\ \sqrt{3})\)，\par\noindent 即斜率\(k\in[\sqrt{3}/3,\ \sqrt{3})\)．}
\end{ansblock}

\tihao{12}\zu{倾斜角与斜率×填空} \tieside{简单(知识点三)}当直线\(l\)的倾斜角\(\theta\in[\pi/4,\pi/2)\cup(\pi/2,2\pi/3]\)时，直线\(l\)的斜率的取值范围为\kongda{[1,+∞)∪(−∞,−√3]}．

\begin{ansblock}[2章-练-课时02-E12]
% ans:2章-练-课时02-E12
\ansitem{12}{[1,+∞)∪(−∞,−√3]}
\ansnote{详解}{\par\noindent \(\theta\in[\pi/4,\pi/2)\)时，\par\noindent \(k=\tan\theta\in[\tan(\pi/4),+\infty)=[1,+\infty)\)；\par\noindent \(\theta\in(\pi/2,2\pi/3]\)时，\par\noindent \(k=\tan\theta\in(-\infty,\tan(2\pi/3)]=(-\infty,-\sqrt{3}]\)．\par\noindent 故\(k\)的取值范围为\([1,+\infty)\cup(-\infty,-\sqrt{3}]\)．}
\end{ansblock}

\tihao{13}\zu{倾斜角与斜率×解答} \tieside{简单(知识点三)}已知直线\(l_1\)的斜率为\(-1\)，直线\(l_2\)的倾斜角比直线\(l_1\)的倾斜角小\(30^\circ\)，求直线\(l_2\)的斜率．
\liubai[16mm]{解答书写区}

\begin{ansblock}[2章-练-课时02-E13]
% ans:2章-练-课时02-E13
\ansitem{13}{−2−√3}
\ansnote{详解}{\(l_1\)的斜率为\(-1\)，其倾斜角为\(135^\circ\)，故\(l_2\)的倾斜角为\(135^\circ-30^\circ=105^\circ\)．\par\noindent \(k_2=\tan 105^\circ=\tan(60^\circ+45^\circ)=(\tan 60^\circ+\tan 45^\circ)/(1-\tan 60^\circ\cdot\tan 45^\circ)=(\sqrt{3}+1)/(1-\sqrt{3})=(\sqrt{3}+1)^{2}/[(1-\sqrt{3})(1+\sqrt{3})]=(4+2\sqrt{3})/(-2)=-2-\sqrt{3}\)．\par\noindent 所以直线\(l_2\)的斜率为\(-2-\sqrt{3}\)．}
\end{ansblock}

\tihao{14}\zu{倾斜角与斜率×解答} \tieside{简单(知识点二)}已知经过\(A(a,-1)\)，\(B(2,a+1)\)的直线的斜率为3，求实数\(a\)的值．
\liubai[16mm]{解答书写区}

\begin{ansblock}[2章-练-课时02-E14]
% ans:2章-练-课时02-E14
\ansitem{14}{a＝1}
\ansnote{详解}{由斜率公式 \(k=[(a+1)-(-1)]/(2-a)=(a+2)/(2-a)=3\)（须\(a\neq 2\)），\par\noindent 去分母得 \(a+2=3(2-a)=6-3a\)，\(4a=4\)，\(a=1\)，且\(a=1\neq 2\)满足条件．所以\(a=1\)．}
\end{ansblock}

% ============================================================
% 思维探索（题15~16＝E15~E16：解答2，居末档照库序）
% ============================================================
\dangbadge{思}{维}{探}{索}

\tihao{15}\zu{倾斜角与斜率×解答} \tieside{简单(知识点一)}已知直线\(l_1\)的倾斜角为\(30^\circ\)，直线\(l_2\)与\(l_1\)垂直，求直线\(l_2\)的倾斜角．
\liubai[16mm]{解答书写区}

\begin{ansblock}[2章-练-课时02-E15]
% ans:2章-练-课时02-E15
\ansitem{15}{120°}
\ansnote{详解}{\(l_2\)与\(l_1\)垂直，则\(l_2\)的向上方向与\(x\)轴正向所成的角为\(30^\circ+90^\circ=120^\circ\)（在\(l_2\)向上方向取与\(l_1\)垂直且指向左上的一侧，另一直线方向与之共线）．\par\noindent 由倾斜角的定义及\(0^\circ\leqslant\alpha<180^\circ\)，得\(l_2\)的倾斜角为\(120^\circ\)．\par\noindent （说明：两直线垂直时倾斜角关系的斜率判据在§2.2.3学习，本题只需用平面几何角的运算作答．）}
\end{ansblock}

\tihao{16}\zu{倾斜角与斜率×解答} \tieside{简单(知识点三)}分别求下列直线的倾斜角：(1) \(y=\sqrt{3}x-2\)；(2) \(x+\sqrt{3}y+1=0\)；(3) \(x=-3\)；(4) \(y=1/2\)．
\liubai[32mm]{解答书写区}

\begin{ansblock}[2章-练-课时02-E16]
% ans:2章-练-课时02-E16
\ansitem{16}{(1) 60°；(2) 150°；(3) 90°；(4) 0°．}
\ansnote{详解}{(1) \(k=\sqrt{3}\)，\(\tan\alpha=\sqrt{3}\)，\(\alpha=60^\circ\)；\par\noindent (2) 化为\(y=-(\sqrt{3}/3)x-\sqrt{3}/3\)，\(k=-\sqrt{3}/3\)，\(\tan\alpha=-\sqrt{3}/3\)，\(\alpha=150^\circ\)；\par\noindent (3) 直线垂直于\(x\)轴，斜率不存在，倾斜角\(\alpha=90^\circ\)；\par\noindent (4) 直线与\(x\)轴平行，倾斜角\(\alpha=0^\circ\)．}
\end{ansblock}

% ---- 尾块（\tailfill 置 \end{multicols} 前；无参＝内容尾随框，每件恰 1 处） ----
\tailfill

\end{multicols}
\end{document}
